\section{Which differences matter?}
The fractions $1/2$ and $2/4$ look different but represent the same rational number. A graph with vertices named Alice, Ben, and Chen may describe the same connection pattern as a graph named $1,2,3$. A necklace does not change when we rotate it, although its written sequence changes if we choose a different starting bead.

In each example, we decide which distinctions to ignore. This decision should be expressed precisely. Otherwise ``these are the same'' can conceal a change to the problem.

A binary relation on a set $A$ specifies which pairs of elements are related. Write $x\sim y$ for one such relation. An \emph{equivalence relation} is a relation with three requirements. \emph{Reflexivity} says each element is related to itself. \emph{Symmetry} says that if $x\sim y$, then $y\sim x$. \emph{Transitivity} says that if $x\sim y$ and $y\sim z$, then $x\sim z$. These are parts of the definition. Calling a new relation an equivalence relation creates three obligations: check every requirement for arbitrary allowed elements, not just a favorite example.

Consider students grouped by their birth month. Each student has the same birth month as themself. If Alice has the same birth month as Ben, Ben has the same birth month as Alice. If Alice and Ben share a birth month and Ben and Chen do too, Alice and Chen share it. The three checks justify regarding students with a common month as one group for this particular comparison. They do not say the students are identical in every other respect.

The relation ``is a friend of'' is generally not transitive, so it is not automatically an equivalence relation. An explanation that says merely ``a relation groups similar things'' is insufficient: similarity may fail the exact conditions needed for a class construction.
